problem:  
Let \((M^n,g,f)\) be a complete noncompact **steady gradient Ricci soliton** satisfying  
\[
\operatorname{Ric} + \nabla^2 f = 0,
\]
with **nonnegative sectional curvature** and all potential level sets compact.  
For each \(r>0\), define the rescaled metric \(g_r = r^{-2} g\) and consider the pointed manifolds \((M,g_r,p)\), where \(p\) is a minimum of \(f\).  
It is conjectured (and known only in special cases such as the Bryant soliton) that such solitons have **asymptotic cones** in the pointed Gromov–Hausdorff sense.

**Open Problem / Conjecture:**  
Without assuming existence of the cone a priori, let \((M,g,f)\) be as above.

1. Does every complete steady gradient Ricci soliton with nonnegative sectional curvature and compact level sets of \(f\) admit, along some sequence \(r_i\to\infty\), a **metric cone limit**
   \[
   (M_\infty, g_\infty,p_\infty) \cong (C(Y), dr^2 + r^2 h),
   \]
   where \(Y\) is a compact length space with \(\operatorname{Ric}_h \ge 0\)?  
2. If such a limit exists, is it necessarily **unique**, i.e., independent of subsequences \(r_i\to\infty\)?  
3. Is the cross‑section \(Y\) always an **Einstein manifold**, or can there exist steady solitons with nonnegative curvature whose blow‑down limits are cones over **non‑Einstein** \(Y\)?  

---

Why is it a "good" problem:  
This problem asks for the *very existence and structure* of asymptotic cones for nonnegatively curved steady Ricci solitons—an issue that lies at the heart of geometric analysis under Ricci flow.

- **Profound insight:**  
  The question probes whether steady solitons always exhibit a well‑defined asymptotic geometric profile. An affirmative answer would generalize key aspects of the Bryant soliton’s rigidity, revealing a universal geometric law for steady singularity models; a counterexample would indicate fundamentally new soliton behaviors. Either outcome would reshape our understanding of Ricci flow singularity formation.  

- **Bridge across fields:**  
  Resolving this problem demands interaction between **analysis of PDEs** (the Ricci soliton equation and curvature estimates), **metric geometry** (Gromov–Hausdorff limits and conical structures), and **Riemannian topology** (properties of positively curved manifolds). It links the fine analytic behavior of the potential \(f\) with the global metric limits of the manifold.  

- **Extreme simplicity:**  
  The question can be summarized in one line:  
  *“Do all nonnegatively curved steady Ricci solitons admit a unique asymptotic cone, and must its cross‑section be Einstein?”*  
  Despite its short and geometrically transparent statement, the issue reaches directly into the uncharted analytic and topological structure of high‑dimensional Ricci solitons—a textbook example of a narrow entrance leading to profound depth.